% concatenated from Korovkin_type_theorems_V4.tex
%\usepackage{showframe}


\documentclass[preprint]{elsarticle}
%%%%%%%%%%%%%%%%%%%%%%%%%%%%%%%%%%%%%%%%%%%%%%%%%%%%%%%%%%%%%%%%%%%%%%%%%%%%%%%%%%%%%%%%%%%%%%%%%%%%%%%%%%%%%%%%%%%%%%%%%%%%%%%%%%%%%%%%%%%%%%%%%%%%%%%%%%%%%%%%%%%%%%%%%%%%%%%%%%%%%%%%%%%%%%%%%%%%%%%%%%%%%%%%%%%%%%%%%%%%%%%%%%%%%%%%%%%%%%%%%%%%%%%%%%%%
\usepackage{amssymb}
\usepackage{amsmath}
\usepackage{amsfonts}
\usepackage{color}
\usepackage{ulem}



\newtheorem{theorem}{Theorem}[section]
\newtheorem{acknowledgement}[theorem]{Acknowledgement}
\newtheorem{algorithm}[theorem]{Algorithm}
\newtheorem{axiom}[theorem]{Axiom}
\newtheorem{case}[theorem]{Case}
\newtheorem{claim}[theorem]{Claim}
\newtheorem{conclusion}[theorem]{Conclusion}
\newtheorem{condition}[theorem]{Condition}
\newtheorem{conjecture}[theorem]{Conjecture}
\newtheorem{corollary}[theorem]{Corollary}
\newtheorem{criterion}[theorem]{Criterion}
\newtheorem{definition}[theorem]{Definition}
\newtheorem{example}[theorem]{Example}
\newtheorem{exercise}[theorem]{Exercise}
\newtheorem{lemma}[theorem]{Lemma}
\newtheorem{notation}[theorem]{Notation}
\newtheorem{problem}[theorem]{Problem}
\newtheorem{proposition}[theorem]{Proposition}
\newtheorem{remark}[theorem]{Remark}
\newtheorem{solution}[theorem]{Solution}
\newtheorem{summary}[theorem]{Summary}
\numberwithin{equation}{section}
\newenvironment{proof}[1][Proof]{\noindent \textbf{#1} }{\  \rule{0.5em}{0.5em}}
\definecolor{c100}{rgb}{0.5,0,.1}
\def \cnm{\color{c100}}
\DeclareMathOperator{\esssup}{esssup}
\journal{...}


\begin{document}
	
	
	\begin{frontmatter}
		
		\title{Korovkin type theorems for operators acting on functions of polynomial and exponential growth on $[0,\infty)$
		}
		
		
		\author[1]{Ulrich Abel} 
		\author[2]{Ana Maria Acu}
		\author[3]{Margareta Heilmann}
		\author[4]{Ioan Ra\c sa}
		
		
		\vspace{10cm}
		
		% \author[]{\hspace{10.5cm} }
		
		
		\address[1] {Technische Hochschule Mittelhessen, Department Mathematik, Naturwissenschaften und Datenverarbeitung, Friedberg, Germany, 
		e-mail: ulrich.abel@mnd.thm.de }
	\address[2]{Lucian Blaga University of Sibiu, Department of Mathematics and Informatics, Sibiu,  Romania, e-mail: anamaria.acu@ulbsibiu.ro}
	\address[3]{University of Wuppertal, School of Mathematics and Natural Sciences, Wuppertal,  Germany, e-mail: heilmann@math.uni-wuppertal.de}
	\address[4]{Technical University of Cluj-Napoca, Department of Mathematics, Cluj-Napoca, Romania,
		e-mail:  ioan.rasa@math.utcluj.ro }
		
		
		
		
		\begin{abstract} 	
			{We prove two Korovkin-type approximation theorems for sequences of positive linear operators acting on continuous functions on $[0,\infty)$. Under the assumption of pointwise convergence on suitable test functions, we establish pointwise convergence for all functions with polynomial or exponential growth.  As direct applications, we obtain convergence results for the classical Baskakov and Szász--Mirakjan operators. The proposed method offers an elementary framework that can be applied to a broad class of positive linear operators.				
			} 
		\end{abstract}
		
		\begin{keyword} Korovkin-type theorem, positive linear operators, polynomial growth, exponential growth
			
			\MSC[2020]  41A36.
		\end{keyword} 
		
	\end{frontmatter}
	
%	\fbox{\fbox{{\large \jobname.tex}}}
	
%	\hskip8.5cm ~ {\large \today}
	
	

\section{Introduction}

Let $L_n : D \to C[0,\infty)$ be a sequence of positive linear operators, where $D$ is a linear subspace of $C[0,\infty)$ containing the space $\Pi$ of all polynomial functions. For example, if $V_n$ is the sequence of Baskakov operators, then $D$ can be taken as the subspace of $C[0,\infty)$ consisting of all functions with at most polynomial growth. In the case of Sz\'asz--Mirakjan operators $S_n$, polynomial growth should be replaced by exponential growth.

For operators $L_n$ as above, many results are known, guaranteeing a certain type of convergence of $L_n f$ toward $f$, for $f$ in suitably restricted subspaces of $D$. In this paper we give sufficient conditions entailing the pointwise convergence of $L_n f$ toward $f$, for {\it all} functions $f \in D$.

\section{Convergence for functions with polynomial growth}

\begin{theorem}\label{T16}Let $D$ be a linear subspace of $C[0,\infty)$, $\Pi \subset D$, $L_n : D \to C[0,\infty)$  positive linear operators, $L_n p(x) \to p(x)$, for all $ p \in \Pi$, $x \geq 0$. Let $f \in D$ such that there exists $ p \in \Pi$ with $|f(x)| \leq p(x)$, $x \geq 0$. Then $L_n f(a) \to f(a)$, $a \geq 0$.\end{theorem}

For the proof we need the following lemma.

\begin{lemma}\label{L6.1} Let $f \in C[0,\infty)$, $p \in \Pi$, $p(x) \geq f(x)$, $x \in [0,\infty)$.
	Let $a \geq 0$, $\varepsilon > 0$. Then, there exists  $P \in \Pi$:
	\begin{equation}\label{equ1}
		P(a) = f(a) + \varepsilon, \quad P(x) \geq f(x), \quad x \geq 0.
	\end{equation}
\end{lemma}
\begin{proof} Let $a>0$. Then, there exists $ 0 < \ell \leq \dfrac{a}{2}$, such that
	\begin{equation}\label{equ2}f(a) + \varepsilon \geq f(x), \,\,x \in (a - 2\ell, a + 2\ell).\end{equation}
	
	
	Let $m \in \mathbb{N}$, $2m > \deg p$. Then, there exists $ M > 0$ with
	\[
	\left( \frac{x - a}{\ell} \right)^{2m} + f(a) + \varepsilon \geq p(x), \quad x > M,
	\]
	which implies that
	\begin{equation}\label{equ3}
		\left( \frac{x - a}{\ell} \right)^{2m} + f(a) + \varepsilon \geq f(x), \quad x > M.
	\end{equation}
	
	Set $U := [0, a - 2\ell] \cup [a + 2\ell, M]$. Take $K \geq 1$ with
	\begin{equation}\label{equ4}
		\min \left\{ K \left( \frac{x - a}{\ell} \right)^{2m} + f(a) + \varepsilon \mid x \in U \right\}
		\geq
		\max \left\{ p(x) \mid x \in U \right\}.
	\end{equation}
	
	Define
	\[
	P(x) := K \left( \frac{x - a}{\ell} \right)^{2m} + f(a) + \varepsilon.
	\]
	
	Using (\ref{equ2})--(\ref{equ4}) we see that $P$ satisfies (\ref{equ1}). The proof is similar for  $a = 0$. \end{proof}


\medskip

\begin{proof}[Proof of Theorem~\protect\ref%
	{T16}]  
	Let $a \geq 0$, $\varepsilon > 0$. Lemma \ref{L6.1} implies the existence of $ P, Q \in \Pi$, such that
	\[
	Q(x) \leq f(x) \leq P(x), \quad x \geq 0; \quad Q(a) = f(a) - \varepsilon,\ \ P(a) = f(a) + \varepsilon.
	\]
	Now
	\[
	L_n Q(a) \leq L_n f(a) \leq L_n P(a);
	\quad f(a) - \varepsilon = Q(a) = \lim_{n\to \infty} L_n Q(a)
	\]
	\[
	\leq \liminf_{n\to\infty} L_n f(a) \leq \limsup_{n\to\infty} L_n f(a)
	\leq \lim_{n\to\infty} L_n P(a) = P(a) = f(a) + \varepsilon. 
	\]
	Therefore, for each $\varepsilon > 0$ we have 
	\[
	f(a) - \varepsilon 
	\leq \liminf_{n\to\infty} L_n f(a) \leq \limsup_{n\to\infty} L_n f(a)
	\leq  f(a) + \varepsilon. 
	\]
	It follows that $L_n f(a) \to f(a)$ as $n \to \infty$.
\end{proof}

\begin{corollary}
	Let $V_n$ be the classical Baskakov operators and $S_n$ the Sz\'asz-Mirakjan operators. Let $f \in C[0,\infty)$ such that there exists $ p \in \Pi$ with $|f(x)| \leq p(x)$, $x \geq 0$.
	Then, $V_n f \to f$ and $S_n f \to f$, pointwise on $[0,\infty)$.\end{corollary}

\begin{proof}
The Baskakov and the Sz\'asz–Mirakjan operators are exponential-type operators (see  \cite[Theorem 3.3]{Ismail}). From \cite[Proposition 2.9]{Ismail} it follows that $V_n p \to p$ and $S_n p \to p$ pointwise on $[0,\infty)$, for all $p \in \Pi$. An application of Theorem \ref{T16} concludes the proof.
	\end{proof}

\section{Convergence for functions with exponential growth}
Let $\exp_{\lambda}(x)=e^{\lambda x}$, $\lambda,x\in{\mathbb R}$. Consider also the function $\alpha(x)=\cosh(x)-1$, $x\in {\mathbb R}$.

\begin{theorem}\label{T5.8} Let \(D\) be a linear subspace of \(C[0,\infty)\) such that \(\exp_\lambda \in D\), \(\lambda \in \mathbb{R}\).
	Let \(L_n : D \to C[0,\infty)\) be positive linear operators, \(L_n \exp_\lambda \to \exp_\lambda\) pointwise on \([0,\infty)\), \(\lambda \in \mathbb{R}\).
	Let \(f \in C[0,\infty)\), \(|f(x)| \le A  \exp_{\omega}(x),\, x\ge 0\) for some \(A>0\) and \(\omega>0\).
	Then \(L_n f \to f\) pointwise on \([0,\infty)\).
\end{theorem}

For the proof we need the following lemma.

	\begin{lemma} Let \(A>0\), \(\omega>0\), \(f \in C[0,\infty)\), such that
	\[
	A \exp_{\omega}(x)  \ge f(x), \quad x \ge 0.
	\]
	Fix \(a \ge 0\), \(\varepsilon > 0\). Then there exist \(\ell>0\), \(K \in \mathbb{N}\), \(m \in \mathbb{N}\) such that
	\begin{equation}
		K\, \alpha^{m}\!\left(\frac{x-a}{\ell}\right) + f(a) + \varepsilon \ge f(x), \quad x \ge 0. \label{ew1}
\end{equation}\end{lemma}

\begin{proof} Let \(a>0\). There exists \(0<\ell \le \dfrac{a}{2}\) such that
\begin{equation}\label{ew2}
	f(a)+\varepsilon \ge f(x), \quad x \in (a-2\ell, a+2\ell). 
\end{equation}

Let \(m \in \mathbb{N}\), \(m > \omega \ell\). Then
\[
\lim_{x \to \infty}
\frac{\alpha^{m}\!\left(\frac{x-a}{\ell}\right) + f(a)+\varepsilon}{A  \exp_{\omega}(x)} = \infty.
\]

So, there exists \(M>0\) such that
\begin{equation}\label{ew3}
	\alpha^{m}\!\left(\frac{x-a}{\ell}\right) + f(a)+\varepsilon \ge f(x), \quad x > M. 
\end{equation}

Set
\[
U := [0, a-2\ell] \cup [a+2\ell, M].
\]
Take \(K \in \mathbb{N}\) with
\begin{equation}\label{ew4}
	\min \left\{ K\, \alpha^{m}\!\left(\frac{x-a}{\ell}\right) + f(a)+\varepsilon \; \middle|\; x \in U \right\}
	\ge
	\max \left\{ A  \exp_{\omega}(x) \; \middle|\; x \in U \right\}. 
\end{equation}

Then (\ref{ew1}) is a consequence of (\ref{ew2})–(\ref{ew4}). The proof is similar for  $a = 0$.
\end{proof}

\medskip

\begin{proof}[Proof of Theorem~\protect\ref%
	{T5.8}]  
	Fix \(a\ge 0\) and \(\varepsilon>0\).
	
	\noindent Since
	\[
	f(x)\le A\exp_\omega(x),\qquad x\ge 0,
	\]
	Lemma 3.2 yields a function
	\[
	G(x)=K\alpha^m\!\left(\frac{x-a}{\ell}\right)+f(a)+\varepsilon
	\]
	such that \(G(x)\ge f(x)\) for all \(x\ge 0\).
	
\noindent	Since \(\alpha(x)=\cosh(x)-1\), the function \(G\) is a finite linear combination of exponentials, hence \(G\in D\). By positivity,
	\[
	L_n f(a)\le L_n G(a).
	\]
	Using the assumption \(\displaystyle\lim_{n\to\infty}L_n\exp_\lambda(a)= \exp_\lambda(a)\) for every \(\lambda\in\mathbb R\), we obtain
	\[
	\displaystyle\lim_{n\to\infty}L_nG(a)= G(a)=f(a)+\varepsilon,
	\]
	and therefore
	\[
	\limsup_{n\to\infty}L_nf(a)\le f(a)+\varepsilon.
	\]
	Applying Lemma 3.2 to \(-f\), we similarly obtain
	\[
	\liminf_{n\to\infty}L_nf(a)\ge f(a)-\varepsilon.
	\]
		Hence
	\[
	f(a)-\varepsilon
	\le
	\liminf_{n\to\infty}L_nf(a)
	\le
	\limsup_{n\to\infty}L_nf(a)
	\le
	f(a)+\varepsilon.
	\]
	Letting \(\varepsilon\to0\), we conclude that
	\[
	\displaystyle\lim_{n\to\infty}L_nf(a)= f(a).
	\]
	Since \(a\ge0\) was arbitrary, the convergence is pointwise on \([0,\infty)\).
	\end{proof}

\begin{corollary}\label{cor5.9} Let \( f \in C[0,\infty) \), \( |f(x)| \le A \exp_{\omega}( x) \), \( x \ge 0 \), for some \( A>0 \) and \( \omega>0 \). Then \( S_n f \to f \) pointwise on \([0,\infty)\).\end{corollary}

\begin{proof} For all $\lambda\in  {\mathbb R}$ and $x\geq 0$, we have
	\[S_n \exp_{\lambda}(x) = e^{n x (e^{\lambda/n}-1)} \to \exp_{\lambda}(x)\quad (n \to \infty).\]
	An application of Theorem \ref{T5.8} concludes the proof. \end{proof}
\begin{thebibliography}{9}
\bibitem{Ismail} M. E. H. Ismail, C. P. May, On a family of approximation operators, J.  Math. Anal. Appl.,  63 (1978), pp. 446–462.
  %\bibitem{2} Zhang, C., Wang, Q. The complete asymptotic expansion for Baskakov operators. Analysis in Theory and Applications 23, 76–82 (2007).
\end{thebibliography}

\end{document}
